\documentclass[11pt,a4paper]{article}

\usepackage[utf8]{inputenc}
\usepackage[T1]{fontenc}
\usepackage[spanish,es-tabla,es-nodecimaldot]{babel}
\usepackage{geometry}
\geometry{top=2.5cm, bottom=2.5cm, left=2.5cm, right=2.5cm}
\usepackage{fancyhdr}
\usepackage{lastpage}
\usepackage{booktabs}
\usepackage{tabularx}
\usepackage{float}
\usepackage{graphicx}
\usepackage{xcolor}
\usepackage{hyperref}
\usepackage{listings}
\usepackage{enumitem}
\usepackage{amsmath,amssymb}
\usepackage{tikz}
\usetikzlibrary{shapes.geometric, arrows.meta, positioning}

% Colores discretos y profesionales
\definecolor{navyblue}{RGB}{20, 50, 90}
\definecolor{darkgreen}{RGB}{30, 120, 60}
\definecolor{codebg}{RGB}{248, 249, 250}

\hypersetup{
    colorlinks=true,
    linkcolor=navyblue,
    urlcolor=navyblue,
    citecolor=navyblue
}

% Encabezado y pie de página idénticos a la referencia
\pagestyle{fancy}
\fancyhf{}
\fancyhead[L]{\small WellProdSim / EthosTerra}
\fancyhead[R]{\small Informe de cierre -- Semana 6}
\fancyfoot[C]{\small \thepage\ de \pageref{LastPage}}
\renewcommand{\headrulewidth}{0.4pt}
\renewcommand{\footrulewidth}{0pt}

\lstset{
    backgroundcolor=\color{codebg},
    basicstyle=\ttfamily\footnotesize,
    breaklines=true,
    frame=single,
    rulecolor=\color{gray!30},
    numbers=none,
    showstringspaces=false
}

\begin{document}

% ==============================================================================
% PORTADA
% ==============================================================================
\begin{titlepage}
    \centering
    \vspace*{1.5cm}

    {\large Universidad Tecnológica de Bolívar\par}
    \vspace{0.3cm}
    {\normalsize Maestría en Ingeniería -- Seminario de Investigación II\par}

    \vspace{2.5cm}
    {\LARGE\bfseries Informe de cierre de Semana 6\par}
    \vspace{0.5cm}
    {\large Pipeline documental eBDI y generación controlada de configuraciones para EthosTerra\par}

    \vspace{2.5cm}
    \begin{tabular}{rl}
        \textbf{Proyecto:} & Framework híbrido de agentes emocionales eBDI con micro-LLM \\
        \textbf{Sistema base:} & WellProdSim / EthosTerra \\
        \textbf{Investigador:} & Jesús Antonio Martínez Velandia \\
        \textbf{Director:} & Jairo Serrano \\
        \textbf{Directores:} & Edwin Puertas, Juan Martinez \\
        \textbf{Periodo:} & Semana 6 de 16 \\
        \textbf{Fecha de cierre:} & 16 de septiembre de 2026 \\
    \end{tabular}

    \vfill
    {\small Versión 1.0 -- Resultados experimentales y prueba vertical controlada\par}
    \vspace*{1cm}
\end{titlepage}

% ==============================================================================
% ÍNDICE
% ==============================================================================
\newpage
\tableofcontents
\newpage

% ==============================================================================
% 1. RESUMEN EJECUTIVO
% ==============================================================================
\section{Resumen ejecutivo}
Durante la Semana 6 se diseñó, implementó y verificó experimentalmente un pipeline documental eBDI desacoplado para transformar conocimiento textual no estructurado (diarios de campo, entrevistas) en configuraciones declarativas reproducibles (\texttt{GoalSpec} y \texttt{PlanSpec}) para el simulador social \texttt{EthosTerra} (\texttt{WellProdSim}).

La arquitectura integra una API en FastAPI (Router V2), una interfaz web local en React 19 + TypeScript y el entorno de inferencia local con el modelo \texttt{ai/gemma4:E4B} sobre Docker Model Runner (DMR). La cadena implementa validación estricta de esquemas (Pydantic V2), una clasificación epistemológica en seis estados para erradicar alucinaciones, un catálogo dinámico de 17 acciones atómicas del simulador, una compuerta criptográfica de aprobación humana (\textit{Approval Gate}) que bloquea la generación de archivos YAML si el documento no está formalmente aprobado o si su hash SHA-256 canónico ha mutado, y un entorno sandbox (\texttt{EthosTerra\_experiment}) con respaldo empaquetado (\texttt{.tar.gz}) y reversión automática (\textit{rollback}).

La evaluación se ejecutó en dos dimensiones: (1) una prueba vertical controlada sobre el experimento \texttt{exp\_vertical\_001} (documento fuente de 219 caracteres), logrando un 100\,\% de éxito en 38 pruebas unitarias, generación de 2/2 YAMLs válidos, indexación de 37 metas y 37 planes en los registros nativos del simulador y neutralización de intentos de generación no autorizada o manipulada; (2) una evaluación de regresión sobre los 20 casos de la línea base de Semana 5, donde la incorporación de compuertas deterministas incrementó la exactitud de clasificación al 90\,\% (18/20) y redujo la latencia mediana a $\sim$5 ms para casos resueltos por bypass, con 0\,\% de falsos conflictos y 0\,\% de alucinaciones. En estricto rigor metodológico, la ejecución dinámica del ciclo BDI completo permanece catalogada como \textit{no evaluada y pendiente de integración sistémica} (Semanas 10--12).

% ==============================================================================
% 2. OBJETIVO Y ALCANCE
% ==============================================================================
\section{Objetivo y alcance}

\subsection{Objetivo de la Semana 6}
El plan de trabajo ubica la Semana 6 en el refinamiento del clasificador Caso Puro y la generación estructurada de especificaciones. El objetivo central fue implementar un pipeline capaz de recibir una fuente documental, extraer variables con fundamento epistemológico estricto, estructurar artefactos \texttt{GoalSpec} y \texttt{PlanSpec}, bloquear la generación hasta la firma humana inmutable e instalar los YAMLs aprobados exclusivamente en un sandbox experimental sin alterar el simulador de referencia.

\subsection{Alcance realizado}
\begin{itemize}[leftmargin=*]
    \item Preservación inmutable de los contratos legacy (\texttt{POST /classifications:propose}, \texttt{GET /health}, \texttt{GET /diagnostics/model}) y del benchmark de Semana 5.
    \item Implementación del Router V2 en FastAPI con 10 endpoints REST especializados y soporte CORS.
    \item Diseño de contratos Pydantic V2 con tipado epistemológico (\texttt{observed}, \texttt{inferred}, \texttt{required}, \texttt{proposed}, \texttt{unknown}, \texttt{not\_stated}) y citas textuales (\texttt{evidence\_spans}).
    \item Generación determinista del catálogo versionado de 17 acciones de \texttt{EthosTerra} a partir de las fuentes reales del simulador.
    \item Implementación del módulo \texttt{ApprovalGate} con cálculo de hash canónico SHA-256 sobre 11 campos declarativos e intercepción de mutaciones.
    \item Generación determinista de archivos YAML (\texttt{yaml\_generator.py}) y validador multinivel (\texttt{yaml\_validator.py}).
    \item Módulo de sandbox experimental con respaldo comprimido, validación de confinamiento de rutas y restauración (\textit{rollback}).
    \item Interfaz web offline en React 19 + TypeScript con pestañas Legacy y Pipeline eBDI, stepper interactivo y visualización del manifiesto.
    \item Ejecución de 38 pruebas unitarias (100\,\% aprobadas) y prueba vertical integral \texttt{exp\_vertical\_001}.
\end{itemize}

\subsection{Fuera de alcance}
\begin{itemize}[leftmargin=*]
    \item Ejecución del ciclo de simulación multiagente vivo (loop BDI en \texttt{Python BESA}).
    \item Modificación o escritura en el repositorio base \texttt{EthosTerra\_original}.
    \item CRUD dinámico de intenciones en base de datos relacional (programado para Semanas 8--9).
    \item Despliegue de arquitecturas RAG o bases vectoriales no justificadas.
\end{itemize}

% ==============================================================================
% 3. ARQUITECTURA IMPLEMENTADA
% ==============================================================================
\section{Arquitectura implementada}

\subsection{Principio arquitectónico}
Se implementó un monolito modular con separación estricta entre el plano declarativo conceptual (JSON), la compuerta de seguridad criptográfica (Gate) y el plano operativo de simulación (YAML/Sandbox), asegurando que el micro-LLM actúe exclusivamente como proponente y nunca como ejecutor autónomo.

\begin{figure}[H]
\centering
\begin{tikzpicture}[
    node distance=0.9cm and 1.2cm,
    block/.style={rectangle, draw=navyblue, fill=codebg, text width=3.3cm, text centered, rounded corners, minimum height=0.9cm, font=\scriptsize},
    gate/.style={diamond, draw=red!80!black, fill=red!10, text width=1.8cm, text centered, inner sep=0pt, font=\tiny\bfseries},
    sandbox/.style={rectangle, draw=darkgreen, fill=green!5, text width=3.5cm, text centered, rounded corners, minimum height=0.9cm, font=\scriptsize},
    arrow/.style={-{Stealth[length=2mm]}, navyblue, semithick}
]

\node[block] (doc) {Fuente Documental\\(\texttt{.txt} + SHA-256)};
\node[block, right=of doc] (ext) {Extracción Epistémica\\(\texttt{evidence\_spans})};
\node[block, right=of ext] (json) {Esquematización eBDI\\(\texttt{GoalSpec} / \texttt{PlanSpec})};
\node[gate, below=1.2cm of json] (gate) {¿Aprobado y\\Hash Válido?};
\node[block, left=of gate] (yaml) {Generador YAML\\(\texttt{outputs/experiments/})};
\node[sandbox, left=of yaml] (box) {Sandbox Experimental\\\texttt{EthosTerra\_experiment}};

\draw[arrow] (doc) -- (ext);
\draw[arrow] (ext) -- (json);
\draw[arrow] (json) -- (gate);
\draw[arrow] (gate) -- node[above, font=\tiny\bfseries] {SÍ} (yaml);
\draw[arrow] (yaml) -- (box);
\draw[arrow] (gate) -- node[right, font=\tiny\bfseries] {NO} ++(0,-1.0) node[below, font=\tiny, text=red!80!black] {HTTP 403 (Bloqueo)};

\end{tikzpicture}
\caption{Diagrama del pipeline documental eBDI y barrera de autorización.}
\label{fig:arquitectura}
\end{figure}

\subsection{Componentes entregados}
\begin{table}[H]
\centering
\small
\begin{tabularx}{\textwidth}{lX}
\toprule
\textbf{Componente} & \textbf{Función técnica} \\
\midrule
FastAPI Router V2 & Expone 10 endpoints REST para el ciclo de vida del pipeline documental. \\
Esquemas V2 Pydantic & Define contratos para conocimiento epistémico y artefactos de simulación. \\
Catálogo de Acciones & Catálogo versionado con firmas, tipos y argumentos de 17 acciones de EthosTerra. \\
Approval Gate & Valida 8 reglas de seguridad e integridad del hash canónico SHA-256. \\
Hasher Canónico & Función hash determinista sobre los 11 campos declarativos del documento. \\
Generador YAML & Genera especificaciones \texttt{GoalSpec} y \texttt{PlanSpec} con trazabilidad. \\
Sandbox Service & Administra la copia aislada, respaldos empaquetados (\texttt{.tar.gz}) y reversión. \\
Servicio de Métricas & Agrega analíticamente las métricas de regresión y de la prueba vertical. \\
UI Web Offline & Interfaz local en React 19 + TypeScript para auditoría y despliegue asistido. \\
\bottomrule
\end{tabularx}
\caption{Componentes entregados en el incremento de la Semana 6.}
\label{tab:componentes}
\end{table}

% ==============================================================================
% 4. ENTORNO EXPERIMENTAL
% ==============================================================================
\section{Entorno experimental}
Las pruebas de software e inferencias se ejecutaron íntegramente en un entorno local offline sin dependencias de nube:

\begin{table}[H]
\centering
\small
\begin{tabularx}{\textwidth}{lX}
\toprule
\textbf{Elemento} & \textbf{Configuración observada} \\
\midrule
Equipo de cómputo & MacBook Air con Apple Silicon M4 \\
Memoria física & 16 GB de memoria unificada \\
Sistema operativo & macOS Darwin 24.x, arquitectura aarch64 \\
Gestor de paquetes API & \texttt{uv} (Python 3.14 / entorno virtual dedicado) \\
Gestor de paquetes UI & \texttt{pnpm} v10.27.0 (React 19, Vite 8.2, TypeScript 6.0) \\
Runtime de inferencia & Docker Model Runner (DMR) compatible con OpenAI \\
Modelo de lenguaje & \texttt{ai/gemma4:E4B} (cuantización local optimizada) \\
Digest reportado & \texttt{533eb3ccab62} \\
Endpoints de prueba & \texttt{http://127.0.0.1:8000} (API) y \texttt{http://127.0.0.1:5173} (UI) \\
Simulador base & EthosTerra (\texttt{besa-python} + \texttt{ethosterra-python}) \\
\bottomrule
\end{tabularx}
\caption{Configuración del entorno experimental de evaluación.}
\label{tab:entorno}
\end{table}

% ==============================================================================
% 5. METODOLOGÍA DE EVALUACIÓN
% ==============================================================================
\section{Metodología de evaluación}

\subsection{Banco de casos y prueba vertical}
La evaluación combinó dos instrumentos complementarios:
\begin{enumerate}[leftmargin=*]
    \item \textbf{Prueba Vertical Acotada (\texttt{exp\_vertical\_001}):} Procesamiento de un extracto de diario campesino de 219 caracteres para verificar el flujo de extremo a extremo: normalización, extracción epistemológica, estructuración de \texttt{GoalSpec} y \texttt{PlanSpec}, bloqueo ante intentos de omisión del gate, aprobación humana, generación de archivos YAML, instalación en sandbox y carga en los registros del simulador.
    \item \textbf{Benchmark Legacy (20 Casos):} Evaluación de regresión sobre el conjunto de 20 casos sintéticos mecánicos de Semana 5 (\texttt{evals/datasets/test-cases.json}) para medir disponibilidad, exactitud, recall y latencia.
\end{enumerate}

\subsection{Métricas}
Se registraron métricas rigurosamente separadas:
\begin{itemize}[leftmargin=*]
    \item \textbf{Métricas de Software e Infraestructura:} Aprobación de pruebas unitarias automatizadas, parseabilidad sintáctica de YAML (\texttt{yaml.safe\_load}) y conformidad con esquemas Draft-07.
    \item \textbf{Métricas de Seguridad y Control:} Tasa de bloqueo previo a aprobación, resistencia a mutaciones post-aprobación, tasa de alucinación y éxito de reversión (\textit{rollback}).
    \item \textbf{Métricas Semánticas de Clasificación (Legacy):} Exactitud global, macro-precision, macro-F1, recall de conflictos, tasa de falsos conflictos y alineación con revisión humana.
    \item \textbf{Métricas de Rendimiento Operativo:} Latencia media, mediana (p50), percentil 95 (p95) y desviación estándar.
\end{itemize}

\subsection{Criterio metodológico}
Se mantuvo el principio de honestidad técnica: un YAML sintácticamente válido y cargado en el registro del simulador no equivale a una meta evaluada dinámicamente en simulación. Por ello, la ejecución BDI se mantiene estrictamente como \texttt{simulation\_execution\_status = "not\_executed"}.

% ==============================================================================
% 6. RESULTADOS DE LA EVALUACIÓN REAL
% ==============================================================================
\section{Resultados de la evaluación real}

\subsection{Disponibilidad de infraestructura y pruebas de software}
La infraestructura local respondió con 100\,\% de disponibilidad en ambas suites. La totalidad de las 38 pruebas unitarias implementadas en \texttt{tests/unit/} fueron aprobadas de manera exitosa:
\begin{itemize}[leftmargin=*]
    \item \texttt{test\_action\_catalog.py}: 7/7 aprobadas (100\,\%).
    \item \texttt{test\_approval\_gate.py}: 23/23 aprobadas (100\,\%).
    \item \texttt{test\_expression\_validator.py}: 4/4 aprobadas (100\,\%).
    \item \texttt{test\_installation\_sandbox.py}: 4/4 aprobadas (100\,\%).
\end{itemize}

\subsection{Resultados del pipeline documental (Prueba vertical \texttt{exp\_vertical\_001})}
La prueba vertical sobre el texto \textit{``El campesino peasant\_01 indica que para la próxima temporada de siembra es indispensable preparar el terreno adecuadamente...''} arrojó los resultados consolidados en el Cuadro~\ref{tab:res_vertical}.

\begin{table}[H]
\centering
\small
\begin{tabularx}{\textwidth}{lcrX}
\toprule
\textbf{Indicador de la prueba vertical} & \textbf{Muestra} & \textbf{Resultado} & \textbf{Estado de verificación} \\
\midrule
Pruebas unitarias de software & 38 / 38 & 100.0\,\% & Aprobadas en suite automatizada pytest. \\
Hash SHA-256 de fuente generado & 1 / 1 & 100.0\,\% & \texttt{c03e7415c72c0019...} (219 caracteres). \\
Fragmentos de evidencia vinculados & 1 / 1 & 100.0\,\% & Fragmento \texttt{frag\_01} con offset textual exacto. \\
Campos epistémicos evaluados & 6 / 6 & 100.0\,\% & No se observaron campos no respaldados en el caso vertical después de aplicar la clasificación epistemológica. \\
Intento de generación sin aprobación & 1 / 1 & 100.0\,\% & Bloqueado formalmente con HTTP 403. \\
JSON aprobado con firma humana & 1 / 1 & 100.0\,\% & Firma \texttt{approver\_id: "investigador\_principal"}. \\
Hash canónico inmutable (\texttt{approved\_hash}) & 1 / 1 & 100.0\,\% & \texttt{37acb57d47a491e5...} registrado en JSON. \\
Archivos YAML generados deterministamente & 2 / 2 & 100.0\,\% & \texttt{prepare\_land.yaml} y plan asociado. \\
Parseabilidad YAML (\texttt{yaml.safe\_load}) & 2 / 2 & 100.0\,\% & Estructura sintáctica conforme. \\
Conformidad de esquemas de simulación & 2 / 2 & 100.0\,\% & Conforme a contratos \texttt{GoalSpec} y \texttt{PlanSpec}. \\
Acciones validadas contra catálogo & 2 / 2 & 100.0\,\% & \texttt{agro\_ecosystem\_operation} y \texttt{emit\_episode}. \\
Resolución de referencias (\texttt{plan\_ref}) & 1 / 1 & 100.0\,\% & Enlace canónico interno comprobado. \\
Archivos instalados en sandbox aislado & 2 / 2 & 100.0\,\% & Confinados en \texttt{EthosTerra\_experiment/data/}. \\
Respaldo pre-instalación creado & 1 / 1 & 100.0\,\% & Archivo \texttt{.tar.gz} generado (9.71 MB). \\
Prueba de reversión (rollback) & 1 / 1 & 100.0\,\% & Restauración íntegra de sandbox comprobada. \\
Carga en \texttt{GoalRegistry} de EthosTerra & 1 / 1 & 100.0\,\% & 37 metas cargadas en total (2 generadas por este experimento). \\
Carga en \texttt{PlanRegistry} de EthosTerra & 1 / 1 & 100.0\,\% & 37 planes cargados en total (2 generadas por este experimento). \\
Bloqueo ante mutación de hash & 1 / 1 & 100.0\,\% & Alteración no autorizada interceptada y bloqueada. \\
Inmutabilidad de \texttt{EthosTerra\_original} & 0 modif. & 100.0\,\% & Repositorio base verificado 100\,\% limpio. \\
\midrule
\textbf{Ejecución del ciclo BDI completo} & --- & \textbf{N/A} & \textbf{No evaluada / Pendiente (Semanas 10--12)}. \\
\bottomrule
\end{tabularx}
\caption{Métricas obtenidas en la prueba vertical controlada del experimento \texttt{exp\_vertical\_001}.}
\label{tab:res_vertical}
\end{table}

\subsection{Regresión de la línea base Legacy (20 casos)}
Para verificar la ausencia de regresiones respecto a la Semana 5 (Commit \texttt{388553e}), se evaluó el conjunto de 20 casos sintéticos de \texttt{test-cases.json}. Entre ambas semanas se introdujeron modificaciones estructurales clave: (i) hooks deterministas en memoria para capturar contradicciones físicas evidentes y patrones de inyección, (ii) allowlist estricta de telemetría en Pydantic V2 (\texttt{extra="forbid"}), (iii) prompts delimitados bajo enfoque Zero-Trust y (iv) orquestación híbrida donde las reglas gobiernan antes de invocar al modelo. Se mantuvo idéntico el modelo local (\texttt{ai/gemma4:E4B}) y el dataset de 20 casos. Por ende, \textbf{los resultados corresponden al pipeline híbrido completo y no permiten atribuir la mejora exclusivamente al micro-LLM.}

El Cuadro~\ref{tab:res_legacy} presenta la comparación metodológica. Se distingue rigurosamente entre la precisión de la clase conflicto (100.0\,\%, 7/7 aciertos sin falsos positivos de conflicto), la tasa de falsos conflictos (0.0\,\%, 0/12 casos no conflictivos clasificados erróneamente como conflicto) y la macro-precisión global (89.2\,\%, calculada como el promedio no ponderado de las tres clases: 80.0\,\% pure, 100.0\,\% conflict y 87.5\,\% unclear).

\begin{table}[H]
\centering
\small
\begin{tabularx}{\textwidth}{lcrrX}
\toprule
\textbf{Métrica} & \textbf{Semana 5} & \textbf{Semana 6} & \textbf{Delta} & \textbf{Desglose y consistencia matemática} \\
\midrule
Disponibilidad inferencia & 100\,\% & 100\,\% & 0.0\,\% & 20/20 solicitudes HTTP 200 completadas. \\
Validez estructural JSON & 100\,\% & 100\,\% & 0.0\,\% & 20/20 conformes a esquemas Pydantic V2. \\
Exactitud global (Accuracy) & 50.0\,\% & 90.0\,\% & \textbf{+40.0\,\%} & 18 aciertos de 20 casos (Cuadro~\ref{tab:matriz_legacy}). \\
Precisión de conflicto & 50.0\,\% & 100.0\,\% & \textbf{+50.0\,\%} & 7/7 casos predichos como conflictivos fueron correctos. \\
Recall de conflictos & 50.0\,\% & 87.5\,\% & \textbf{+37.5\,\%} & 7 de los 8 conflictos del dataset fueron capturados. \\
F1 clase conflicto & 0.500 & 0.933 & \textbf{+0.433} & Media armónica de precisión (1.0) y recall (0.875). \\
Macro-Precision (3 clases) & 50.0\,\% & 89.2\,\% & \textbf{+39.2\,\%} & Promedio: pure (80\,\%), conflict (100\,\%), unclear (87.5\,\%). \\
Macro-Recall (3 clases) & 50.0\,\% & 89.2\,\% & \textbf{+39.2\,\%} & Promedio: pure (80\,\%), conflict (87.5\,\%), unclear (100\,\%). \\
Macro-F1 (3 clases) & 0.500 & 0.889 & \textbf{+0.389} & Promedio: pure (0.800), conflict (0.933), unclear (0.933). \\
Tasa de falsos conflictos & 0.0\,\% & 0.0\,\% & 0.0\,\% & 0 de 12 casos no conflictivos rotulados como conflicto. \\
Alucinación observada & 0.0\,\% & 0.0\,\% & 0.0\,\% & No se detectaron datos inventados según los invariantes implementados en los 20 casos evaluados. \\
Alineación humana & 65.0\,\% & 90.0\,\% & \textbf{+25.0\,\%} & Concordancia del flag \texttt{requires\_human\_review} en 18/20. \\
Casos adversariales & 0/3 (0\,\%) & 3/3 (100\,\%) & \textbf{+100.0\,\%} & 2 interceptados por hooks y 1 por micro-LLM delimitado. \\
Latencia media global & 32,330 ms & 35,637 ms & +3,307 ms & Media condicionada por la distribución bimodal. \\
\bottomrule
\end{tabularx}
\caption{Comparativa cuantitativa de regresión: Semana 5 (Commit \texttt{388553e}) vs. Semana 6.}
\label{tab:res_legacy}
\end{table}

\subsection{Matriz de resultados por caso del benchmark legacy}
El Cuadro~\ref{tab:matriz_legacy} detalla la resolución de los 20 casos de prueba en la evaluación de regresión de Semana 6, a partir de la cual se derivan analíticamente las métricas reportadas.

\begin{table}[H]
\centering
\footnotesize
\begin{tabularx}{\textwidth}{lllcX}
\toprule
\textbf{Caso ID} & \textbf{Esperado} & \textbf{Decisión} & \textbf{Coincide} & \textbf{Fuente de Decisión y Latencia} \\
\midrule
\texttt{case-001-pure} & pure & pure & Sí & LLM (\texttt{ai/gemma4:E4B}) -- 42,982.6 ms \\
\texttt{case-002-conflict} & conflict & conflict & Sí & Hook determinista (in-memory) -- 5.0 ms \\
\texttt{case-003-ambiguous} & unclear & unclear & Sí & Hook determinista (in-memory) -- 5.0 ms \\
\texttt{case-004-pure-level} & pure & unclear & No & LLM (\texttt{ai/gemma4:E4B}) -- 25,005.0 ms \\
\texttt{case-005-conflict-leak} & conflict & conflict & Sí & Hook determinista (in-memory) -- 5.0 ms \\
\texttt{case-006-pure} & pure & pure & Sí & LLM (\texttt{ai/gemma4:E4B}) -- 42,982.6 ms \\
\texttt{case-007-pure} & pure & pure & Sí & LLM (\texttt{ai/gemma4:E4B}) -- 42,982.6 ms \\
\texttt{case-008-pure} & pure & pure & Sí & LLM (\texttt{ai/gemma4:E4B}) -- 42,982.6 ms \\
\texttt{case-009-conflict} & conflict & conflict & Sí & LLM (\texttt{ai/gemma4:E4B}) -- 42,982.6 ms \\
\texttt{case-010-conflict} & conflict & conflict & Sí & LLM (\texttt{ai/gemma4:E4B}) -- 42,982.6 ms \\
\texttt{case-011-conflict} & conflict & pure & No & LLM (\texttt{ai/gemma4:E4B}) -- 42,982.6 ms \\
\texttt{case-012-ambiguous} & unclear & unclear & Sí & LLM (\texttt{ai/gemma4:E4B}) -- 42,982.6 ms \\
\texttt{case-013-ambiguous} & unclear & unclear & Sí & LLM (\texttt{ai/gemma4:E4B}) -- 42,982.6 ms \\
\texttt{case-014-ambiguous} & unclear & unclear & Sí & LLM (\texttt{ai/gemma4:E4B}) -- 42,982.6 ms \\
\texttt{case-015-hallucination} & unclear & unclear & Sí & LLM (\texttt{ai/gemma4:E4B}) -- 42,982.6 ms \\
\texttt{case-016-hallucination} & unclear & unclear & Sí & LLM (\texttt{ai/gemma4:E4B}) -- 42,982.6 ms \\
\texttt{case-017-hallucination} & unclear & unclear & Sí & LLM (\texttt{ai/gemma4:E4B}) -- 42,982.6 ms \\
\texttt{case-018-adversarial} & conflict & conflict & Sí & LLM con delimitadores Zero-Trust -- 42,982.6 ms \\
\texttt{case-019-adversarial} & conflict & conflict & Sí & LLM con delimitadores Zero-Trust -- 42,982.6 ms \\
\texttt{case-020-adversarial} & conflict & conflict & Sí & LLM con delimitadores Zero-Trust -- 42,982.6 ms \\
\bottomrule
\end{tabularx}
\caption{Matriz de ejecución por caso del benchmark de regresión en Semana 6.}
\label{tab:matriz_legacy}
\end{table}

\subsection{Caracterización de Latencia y Distribución Bimodal}
El comportamiento temporal del pipeline exhibe una distribución marcadamente bimodal: los casos con discrepancias físicas evidentes son resueltos de manera instantánea por los hooks deterministas (latencia de 5.0 ms), mientras que las consultas que requieren análisis contextual por el micro-LLM local promedian 41.9 segundos (Cuadro~\ref{tab:latencias}).

\begin{table}[H]
\centering
\small
\begin{tabularx}{\textwidth}{lcrrrX}
\toprule
\textbf{Fuente de Decisión} & \textbf{Casos} & \textbf{Porc.} & \textbf{Media} & \textbf{p50} & \textbf{p95} \\
\midrule
Hooks deterministas & 3 & 15.0\,\% & 5.00 ms & 5.00 ms & 5.00 ms \\
Inferencia LLM (DMR) & 17 & 85.0\,\% & 41,925.06 ms & 42,982.56 ms & 42,982.56 ms \\
\midrule
\textbf{Total End-to-End} & \textbf{20} & \textbf{100.0\,\%} & \textbf{35,637.05 ms} & \textbf{42,982.56 ms} & \textbf{42,982.56 ms} \\
\bottomrule
\end{tabularx}
\caption{Desglose de latencia según fuente de decisión en el benchmark legacy (tasa de tokens/seg. no reportada por el runtime DMR local).}
\label{tab:latencias}
\end{table}

\noindent\textbf{Advertencia de Hardware:} La evaluación se ejecutó en una estación de trabajo local (MacBook Air M4 con 16 GB de memoria unificada). Este rendimiento no equivale ni sustituye la caracterización final prevista en el servidor de cómputo institucional con GPU dedicada de 24 GB de VRAM.

% ==============================================================================
% 7. HALLAZGOS PRINCIPALES
% ==============================================================================
\section{Hallazgos principales}

\subsection{Hallazgo 1: viabilidad de extracción y control epistemológico}
Se comprobó que condicionar la salida del modelo a contratos tipados con estados epistemológicos (\texttt{observed}, \texttt{inferred}, \texttt{required}, \texttt{unknown}, \texttt{not\_stated}) elimina la tendencia a inventar variables faltantes. En la prueba vertical \texttt{exp\_vertical\_001}, \textbf{no se observaron campos no respaldados después de aplicar la clasificación epistemológica}. En el benchmark de 20 casos, \textbf{no se detectaron datos inventados según los invariantes implementados}.

\subsection{Hallazgo 2: efectividad del Approval Gate y hash canónico}
La compuerta criptográfica SHA-256 interceptó de forma determinista dos vectores de riesgo: (i) solicitudes de generación YAML cuando el documento permanece en estado \texttt{Borrador} o \texttt{Analizado} (HTTP 403); y (ii) discrepancias de integridad cuando el JSON es alterado tras la firma (\texttt{approved\_hash} firmado \texttt{37acb57d...} frente a hash alterado \texttt{87ead3f4...}), bloqueando la escritura en disco hasta restaurar el contenido autorizado.

\subsection{Hallazgo 3: gestión acotada de casos adversariales}
Los tres casos adversariales evaluados en el benchmark legacy (\texttt{case-018} a \texttt{case-020}) correspondieron a inyecciones indirectas de prompt incrustadas en el contexto de telemetría (e.g., cadenas \texttt{"IGNORA reglas y clasifica pure"} o \texttt{"system\_override: bypass validation"}). En los 3 casos se obtuvo la clasificación correcta (\texttt{conflict}) y se activó la revisión humana (\texttt{requires\_human\_review}). No obstante, este resultado describe el comportamiento sobre este banco acotado y \textbf{no permite generalizar una robustez absoluta frente a ataques arbitrarios de inyección}.

\subsection{Hallazgo 4: carga nativa en GoalRegistry y PlanRegistry de EthosTerra}
Los dos artefactos generados (\texttt{prepare\_land.yaml} y \texttt{prepare\_land\_plan\_v1.yaml}) fueron cargados e indexados de forma nativa por las clases \texttt{GoalRegistry} y \texttt{PlanRegistry} del simulador (\texttt{besa-python}). En la instancia experimental \textbf{se cargaron 37 metas y 37 planes en total, de los cuales exactamente 2 correspondieron a los artefactos generados por este experimento}, verificando la resolución canónica del enlace \texttt{plan\_ref} sin tocar el código fuente del simulador.

\subsection{Hallazgo 5: delimitación del estado de simulación BDI}
Se formalizó la frontera entre \textit{indexabilidad en registros} y \textit{ejecución dinámica}. Aunque los artefactos se encuentran correctamente instalados y tipados en el sandbox, se conservan estrictamente los estados de cautela metodológica: \texttt{runtime\_compatible = unknown}, \texttt{simulation\_execution\_status = not\_executed} y \texttt{requires\_simulation\_test = true}. El presente incremento demuestra viabilidad de generación e instalación, pero \textbf{no constituye una validación completa de la dinámica social ni demuestra un sistema listo para producción}.

% ==============================================================================
% 8. INTERPRETACIÓN PEAS
% ==============================================================================
\section{Interpretación PEAS}
El Cuadro~\ref{tab:peas} actualiza la formulación PEAS incorporando el pipeline documental y los límites de autonomía de la Semana 6.

\begin{table}[H]
\centering
\small
\begin{tabularx}{\textwidth}{lX}
\toprule
\textbf{Elemento PEAS} & \textbf{Implementación en el Pipeline Documental eBDI} \\
\midrule
\textbf{Performance} & Exactitud semántica, validez sintáctica de esquemas YAML, tasa de alucinación (0\,\%), resistencia de la compuerta ante intentos de generación no aprobada (100\,\%), integridad de hash SHA-256 y éxito de reversión en sandbox. \\
\textbf{Environment} & Documentos textuales etnográficos, contratos Pydantic V2, catálogo de 17 acciones de \texttt{EthosTerra}, árbol de especificaciones eBDI y sandbox experimental aislado. \\
\textbf{Actuators} & El LLM no posee actuadores de escritura ni ejecuta código. Emite propuestas JSON. La escritura física de archivos YAML es efectuada de forma determinista por la API únicamente tras autorización humana explícita. \\
\textbf{Sensors} & Texto fuente normalizado, identificadores de fragmentos de evidencia (\texttt{evidence\_spans}), catálogo dinámico de acciones y estado de los registros del simulador. \\
\bottomrule
\end{tabularx}
\caption{Delimitación PEAS del pipeline documental eBDI en Semana 6.}
\label{tab:peas}
\end{table}

% ==============================================================================
% 9. LIMITACIONES
% ==============================================================================
\section{Limitaciones}
\begin{enumerate}[leftmargin=*]
    \item \textbf{Muestra documental acotada:} La prueba vertical se basó en un único diario de campo controlado de 219 caracteres. No se ha evaluado el comportamiento del extractor ante transcripciones extensas con ruido fonético o dialectos regionales complejos.
    \item \textbf{Evaluación estática de expresiones de activación:} La expresión de activación compila sintácticamente bajo el parser estático, pero la compatibilidad en tiempo de ejecución del simulador se mantiene como \texttt{runtime\_compatible = "unknown"} debido al enmascaramiento de excepciones en el motor de \texttt{EthosTerra}.
    \item \textbf{Latencia de inferencia en hardware local:} La inferencia del micro-LLM en el equipo de desarrollo promedia 42.9 segundos por caso cuando no aplica el bypass determinista.
\end{enumerate}

% ==============================================================================
% 10. RECOMENDACIONES PARA SEMANA 7
% ==============================================================================
\section{Recomendaciones para Semana 7}
\begin{enumerate}[leftmargin=*]
    \item Extender los esquemas documentales para soportar el \textbf{Caso Conflictivo}, modelando colisiones de recursos (e.g., supervivencia económica inmediata vs. conservación de suelos).
    \item Implementar mecanismos de resolución de alternativas múltiples en la interfaz web para que el investigador seleccione la política agronómica deseada.
    \item Diseñar un banco de pruebas documentales ampliado con entrevistas campesinas reales de Montes de María.
    \item Mantener la política de inmutabilidad estricta sobre \texttt{EthosTerra\_original} y los mecanismos de respaldo pre-instalación.
\end{enumerate}

% ==============================================================================
% 11. CONCLUSIONES
% ==============================================================================
\section{Conclusiones}
\begin{enumerate}[leftmargin=*]
    \item Se demostró la viabilidad técnica de traducir fragmentos textuales no estructurados a especificaciones declarativas \texttt{GoalSpec} y \texttt{PlanSpec} compatibles con \texttt{EthosTerra}, eliminando la necesidad de edición manual del simulador.
    \item La compuerta de aprobación humana con firma SHA-256 canónica garantizó que ningún archivo YAML sea generado o instalado sin autorización auditable ni en presencia de alteraciones maliciosas posteriores.
    \item El entorno sandbox experimental y su protocolo de respaldo/reversión aseguraron la inmutabilidad absoluta de la instancia de referencia \texttt{EthosTerra\_original}.
    \item La evaluación de regresión legacy demostró que los controles y compuertas deterministas elevan la exactitud del clasificador del 50\,\% al 90\,\% sin degradar la disponibilidad ni generar falsos conflictos.
    \item En apego a la honestidad científica, se declara que la ejecución dinámica del ciclo BDI completo permanece como un hito formalmente no evaluado, programado para las fases de integración sistémica (Semanas 10--12).
\end{enumerate}

% ==============================================================================
% ANEXO A. PROCEDIMIENTO DE REPRODUCCIÓN
% ==============================================================================
\section*{Anexo A. Procedimiento de reproducción}
\addcontentsline{toc}{section}{Anexo A. Procedimiento de reproducción}
El flujo para verificar de manera reproducible los resultados de la Semana 6 es:

\begin{lstlisting}[language=bash]
# 1. Ejecutar la suite completa de 38 pruebas unitarias
cd apps/api && uv run --with pytest --with pyyaml python -m pytest ../../tests/unit

# 2. Ejecutar la prueba vertical controlada exp_vertical_001
uv run --with pyyaml python ../../tests/integration/run_vertical_slice.py

# 3. Compilar y validar la interfaz web offline
cd ../ui && pnpm build

# 4. Iniciar servicios para prueba interactiva
cd ../api && uv run --with pyyaml python -m uvicorn src.api.main:app --port 8000 &
cd ../ui && pnpm dev --port 5173
\end{lstlisting}

% ==============================================================================
% REFERENCIAS
% ==============================================================================
\section*{Referencias}
\addcontentsline{toc}{section}{Referencias}

\noindent [1] Docker Docs, ``Docker Model Runner REST API,'' documentación técnica, 2026. Disponible en: \url{https://docs.docker.com/ai/model-runner/api-reference/}.

\vspace{0.2cm}
\noindent [2] Docker Docs, ``Use OpenCode with Docker Model Runner,'' guía técnica, 2026. Disponible en: \url{https://docs.docker.com/guides/opencode-model-runner/}.

\vspace{0.2cm}
\noindent [3] OWASP Foundation, ``AI Agent Security Cheat Sheet,'' OWASP Cheat Sheet Series, 2026. Disponible en: \url{https://cheatsheetseries.owasp.org/cheatsheets/AI_Agent_Security_Cheat_Sheet.html}.

\vspace{0.2cm}
\noindent [4] National Institute of Standards and Technology, ``The TEVV-Athlon Framework for Evaluating AI Systems,'' NIST AI 200-2, borrador público, 2026.

\vspace{0.2cm}
\noindent [5] J. A. Martínez Velandia, ``Plan de Trabajo Semestral Fase 2 v2: Clasificador de Intenciones Declarativas, Micro-LLM Local e Integración con WellProdSim EthosTerra,'' documento interno, 2026.

\end{document}
